\begin{thebibliography}{99}

\bibitem{bean2025} Bean, A. M., Kearns, R. O., Romanou, A., Hafner, F. S., Mayne, H., et al. (2025). Measuring what Matters: Construct Validity in Large Language Model Benchmarks. NeurIPS 2025 Datasets and Benchmarks Track.

\bibitem{kearns2026} Kearns, R. O. (2026). Quantifying construct validity in large language model evaluations. MSc thesis, Oxford Internet Institute, University of Oxford. arXiv:2602.15532.

\bibitem{akhtar2026} Akhtar, M., Reuel, A., Soni, P., Ahuja, S., et al. (2026). When AI Benchmarks Plateau: A Systematic Study of Benchmark Saturation. Preprint, Evaluating Evaluations (EvalEval) Coalition. arXiv:2602.16763.

\bibitem{reuel2026} Reuel, A., Ghosh, A., Chim, J., Tran, A., et al. (2026). Who Evaluates AI's Social Impacts? Mapping Coverage and Gaps in First and Third Party Evaluations. Preprint, Evaluating Evaluations (EvalEval) Coalition. arXiv:2511.05613.

\bibitem{thurstone1947} Thurstone, L. L. (1947). Multiple-Factor Analysis. University of Chicago Press.

\bibitem{horn1965} Horn, J. L. (1965). A rationale and test for the number of factors in factor analysis. Psychometrika, 30(2), 179-185.

\bibitem{cattell1966} Cattell, R. B. (1966). The scree test for the number of factors. Multivariate Behavioral Research, 1(2), 245-276.

\bibitem{kaiser1974} Kaiser, H. F. (1974). An index of factorial simplicity. Psychometrika, 39(1), 31-36.

\bibitem{bartlett1950} Bartlett, M. S. (1950). Tests of significance in factor analysis. British Journal of Psychology (Statistical Section), 3, 77-85.

\bibitem{rousseeuw1987} Rousseeuw, P. J. (1987). Silhouettes: a graphical aid to the interpretation and validation of cluster analysis. Journal of Computational and Applied Mathematics, 20, 53-65.

\bibitem{tibshirani2001gap} Tibshirani, R., Walther, G., \& Hastie, T. (2001). Estimating the number of clusters in a data set via the gap statistic. Journal of the Royal Statistical Society: Series B, 63(2), 411-423.

\bibitem{hoerl1970} Hoerl, A. E., \& Kennard, R. W. (1970). Ridge regression: biased estimation for nonorthogonal problems. Technometrics, 12(1), 55-67.

\bibitem{zou2005} Zou, H., \& Hastie, T. (2005). Regularization and variable selection via the elastic net. Journal of the Royal Statistical Society: Series B, 67(2), 301-320.

\bibitem{breiman2001} Breiman, L. (2001). Random forests. Machine Learning, 45(1), 5-32.

\bibitem{friedman2001} Friedman, J. H. (2001). Greedy function approximation: a gradient boosting machine. Annals of Statistics, 29(5), 1189-1232.

\bibitem{lundberg2017} Lundberg, S. M., \& Lee, S.-I. (2017). A unified approach to interpreting model predictions. Advances in Neural Information Processing Systems 30 (NeurIPS 2017), 4765-4774.

\bibitem{efron1993} Efron, B., \& Tibshirani, R. J. (1993). An Introduction to the Bootstrap. Chapman \& Hall.

\bibitem{pedregosa2011} Pedregosa, F., Varoquaux, G., Gramfort, A., Michel, V., et al. (2011). Scikit-learn: Machine Learning in Python. Journal of Machine Learning Research, 12, 2825-2830.

\bibitem{aa2026} Artificial Analysis. (2026). Artificial Analysis LLM Leaderboard and Intelligence Index. https://artificialanalysis.ai (snapshot accessed 6 July 2026).

\bibitem{epoch2026} Epoch AI. (2026). Notable AI Models (data set, CC-BY 4.0). https://epoch.ai/data/notable-ai-models.

\bibitem{hastie2009} Hastie, T., Tibshirani, R., \& Friedman, J. (2009). The Elements of Statistical Learning: Data Mining, Inference, and Prediction (2nd ed.). Springer.
\end{thebibliography}
